% preamble.tex

% --- Pacchetti Fondamentali ---
\usepackage[utf8]{inputenc}
\usepackage[T1]{fontenc}
\usepackage[english]{babel}
\usepackage{csquotes}
\usepackage{amsmath, amssymb, amsthm} % Ambienti matematici
\usepackage{mathrsfs} % Per il corsivo calligrafico (\mathscr)
\usepackage{mathtools} % Estensioni per amsmath
\usepackage{geometry}
\geometry{a4paper, top=3cm, bottom=3cm, left=3.5cm, right=2.5cm}
\usepackage{microtype}

% --- Gestione dei Subfiles ---
\usepackage{subfiles}

% --- Bibliografia (con biblatex) ---
\usepackage[backend=biber, style=authoryear, sorting=nyt]{biblatex}
\addbibresource{references.bib}

% --- Collegamenti Ipertestuali ---
\usepackage[colorlinks=true, linkcolor=blue, citecolor=blue, urlcolor=blue]{hyperref}

% --- Ambiente per Teoremi e Dimostrazioni ---
% Stile "plain" per i risultati principali
\theoremstyle{plain}
\newtheorem{theorem}{Teorema}[chapter]
\newtheorem{lemma}[theorem]{Lemma}
\newtheorem{corollary}[theorem]{Corollario}
\newtheorem{proposition}[theorem]{Proposizione}

% Stile "definition" per definizioni ed esempi
\theoremstyle{definition}
\newtheorem{definition}[theorem]{Definizione}
\newtheorem{example}[theorem]{Esempio}
\newtheorem{exercise}[theorem]{Esercizio}

% Stile "remark" per osservazioni e note (non numerati)
\theoremstyle{remark}
\newtheorem*{remark}{Osservazione}
\newtheorem*{notation}{Notazione}
\newtheorem*{abuse}{Abuso di notazione}

% comandi personalizzati
\newcommand{\J}{\mathcal{J}}
\newcommand{\JJ}{\mathcal{\overline{J}}}
\newcommand{\U}{\mathcal{U}}
\newcommand{\R}{\mathcal{R}}
\newcommand{\HSS}{Hom(S_x(M),S_x(M))}
\newcommand{\SxM}{S_x(M)}
\newcommand{\Lang}{\mathcal{L}}
\newcommand{\LLang}{\overline{\Lang}}
\newcommand{\nonfork}{%
  \mathop{\raise0.2ex\hbox{%
    \ooalign{\hidewidth$\vert$\hidewidth\cr
             \raise-0.9ex\hbox{$\smile$}}}}%
}

